\section{Power Divergences and Renormalization of Quasi-Distributions in Lattice QCD}

In a paper by Rossi and Resta~\cite{Rossi:2017muf}, a question has been raised about the practical implementation
of the factorization theorem, Eq.~(\ref{eq:fact}), in lattice QCD. In particular, they consider
the quasi-distribution in Eq. (19) of their paper, and try to
compute its second moment in the rest frame of the hadron. They point out that
because of the mixing of the trace term with the lower dimensional operator, the
$1/p_z^2$ suppressed terms have a power-divergent contribution as shown in Eq. (35) of their paper.
Therefore, they conclude that ``unless one performs a non-perturbative subtraction of
the power divergent terms, the $p_z\rightarrow \infty$ limit does not exist.''
However, a careful analysis shows that this problem does not exist in the present formalism.

First of all, the quasi-distribution used to extract the physical
parton distribution, is a function of $y=k_z/p_z$. Formally, one can take the second
moment and show that it results in a local operator with double derivatives.
However, in practice, the moments of the quasi-distribution do not exist. As can be seen from the one-loop calculation~\cite{Xiong:2013bka},
the result behaves like $1/|y|$ at large $y$, therefore all higher
moments other than the zeroth one do not converge.
It is exactly this large $y$ divergence that calls for an extra renormalization
in the local operators. However, without taking the moments, the divergence
does not appear.

In fact, it is exactly this non-local quasi-distribution formulation of parton physics that
avoids the power divergence problem in moment calculations using the old-fashioned
approach. The non-local formulation has much simpler ultraviolet (UV) physics
and hence makes it much easier to control the divergences.

For example, the renormalization of the quasi-distributions is straightforward in ``heavy-quark'' formalism. It can be shown
that for quark non-singlets, except for the Wilson-line self-energy or mass correction,
there are no other power divergences~\cite{Ishikawa:2016znu,Chen:2016fxx,Ji:2017oey,Ishikawa:2017faj,Green:2017xeu}. All other divergences are logarithmic
in the lattice spacing. Therefore, the quasi-distributions can be
renormalized by factoring out the power-divergent self-energy factor and logarithmic
renormalization factor, and have a smooth limit when the lattice spacing vanishes.
Note that without factoring out the power-divergent self-energy factor,
a one-loop mass correction of form $\alpha_s C_1|z|/a$ with lattice spacing $a$
and Wilson-line length $|z|$ can generate a finite collinearly divergent term
in the quasi-distribution at two-loop, which appears to break the IR factorization
of the quasi-distribution~\cite{Li:2016amo}. However, such terms do not
contribute when the linear divergence is factored out of the quasi-distributions before taking
the lattice spacing to zero.

After renormalizing the quasi-distributions in an appropriate scheme, one can then use the factorization formula in Eq.~(\ref{eq:fact}) to match it to the $\overline{\text{MS}}$ parton distribution. Then the higher-twist $1/p_z^2$ terms vanish smoothly as $p_z\rightarrow \infty$.
This is exactly what RI/MOM scheme does in recent papers~\cite{Alexandrou:2017huk,Chen:2017mzz,Stewart:2017tvs}.
